\subsection{Cone of a map}

Let X and Y be topological spaces, and let f be a continuous map from X to Y. Let \(\operatorname{Cyl}(f)\) be the cylinder of the map f. Let us denote \(\alpha_{f}\) the canonical map from \(X \times I\) into \(\operatorname{Cyl}(f)\) and \(f_{0}\) the map \(x \mapsto \alpha_{f}(x,0)\) from X to \(\operatorname{Cyl}(f)\); these maps induce homeomorphisms from \(X \times I\) and X respectively onto their images in \(\operatorname{Cyl}(f)\).

DEFINITION 11. --- The cone of the map f is denoted by \(\operatorname{Côn}(f)\) and is the topological space obtained from \(\operatorname{Cyl}(f)\) by contracting \(f_{0}(\mathrm{X})\).

Denote by \(\beta_{f}^{\prime}\colon\mathrm{Y}\to\mathrm{C}\hat{\mathrm{o}}\mathrm{n}(f)\) the composition of the canonical map \(\beta_{f}\colon\mathrm{Y}\to\mathrm{C}\mathrm{y}\mathrm{l}(f)\) and the canonical surjection from \(\mathrm{C}\mathrm{y}\mathrm{l}(f)\) onto \(\mathrm{C}\hat{\mathrm{o}}\mathrm{n}(f)\). The map \(\beta_{f}^{\prime}\) is continuous and defines a homeomorphism from Y onto a closed subset of \(\mathrm{C}\hat{\mathrm{o}}\mathrm{n}(f)\), called the base of the cone, which we shall thereby identify with Y.

Denote by \(\alpha_{f}^{\prime}\colon\mathrm{X}\times\mathbf{I}\to\mathrm{C}\hat{\mathrm{o}}\mathrm{n}(f)\) the composition of the map \(\alpha_{f}\) and the canonical surjection from \(\mathrm{Cyl}(f)\) onto \(\mathrm{C}\hat{\mathrm{o}}\mathrm{n}(f)\) . The map \(\alpha_{f}^{\prime}\) is a homotopy whose initial map is a constant map and whose terminal map is the map \(\beta_{f}^{\prime}\circ f\) .

If X is empty, the map \(\beta_{f}^{\prime}\) is a homeomorphism from Y onto \(\operatorname{Côn}(f)\) .

Suppose that X is not empty. Let s denote the base point of the space \(\mathrm{Cyl}(f)/f_{0}(\mathrm{X})\); it is called the vertex of the cone \(\mathrm{Côn}(f)\). Since \(f_{0}(\mathrm{X})\) is closed in \(\mathrm{Cyl}(f)\), the canonical map \(\pi\colon\mathrm{Cyl}(f)\to\mathrm{Côn}(f)\) induces a homeomorphism of \(\mathrm{Cyl}(f)\setminus f_{0}(\mathrm{X})\) onto \(\mathrm{Côn}(f)\setminus\{s\}\) (III, p. 252). The vertex of the cone \(\mathrm{Côn}(f)\) does not belong to its base; the canonical injection of Y into \(\mathrm{Côn}(f)\setminus\{s\}\) is a homotopy equivalence (III, p. 239, prop. 6).

Let \(\sigma_f\colon \mathrm{Cyl}(f)\times \mathbf{I}\to \mathrm{Cyl}(f)\) be the canonical contraction of the cylinder of \(f\) onto its base. For \(c\in \mathrm{Cyl}(f)\setminus f_0(\mathrm{X})\) and \(t\in \mathbf{I}\), we have \(\sigma_f(c,t)\notin f_0(\mathrm{X})\). Consequently, there exists a unique map \(\sigma_f'\colon (\mathrm{Côn}(f)\setminus\{s\})\times \mathbf{I}\to \mathrm{Côn}(f)\setminus\{s\}\) such that \(\sigma_f'(\pi(c),t)=\pi(\sigma_f(c,t))\) for \(c\in \mathrm{Cyl}(f)\setminus f_0(\mathrm{X})\) and \(t\in \mathbf{I}\). It is continuous and is a strong contraction of \(\mathrm{Côn}(f)\setminus\{s\}\) onto Y. It is called the canonical contraction of \(\mathrm{Côn}(f)\setminus\{s\}\) onto its base. Its terminal map is a retraction of \(\mathrm{Côn}(f)\setminus\{s\}\) onto Y, called the canonical retraction of the cone with its vertex removed onto its base.

PROPOSITION 11 (Universal property of cones)

Let Z be a topological space, let \(\beta: \mathrm{Y} \to \mathrm{Z}\) be a continuous map and let \(\alpha: \mathrm{X} \times \mathbf{I} \to \mathrm{Z}\) be a homotopy whose initial map is a constant map and whose terminal map is equal to \(\beta \circ f\). There exists a unique continuous map \(\varphi\) from \(\operatorname{Côn}(f)\) into Z such that \(\alpha = \varphi \circ \alpha_f'\) and \(\beta = \varphi \circ \beta_f'\).

By the universal property of cylinders (III, p. 238, prop. 4), there exists a unique continuous map \(h \colon \operatorname{Cyl}(f) \to \mathbf{Z}\) such that \(\alpha = h \circ \alpha_f\) and \(\beta = h \circ \beta_f\) . Since the initial map of \(\alpha\) is a constant map, the restriction of \(h\) to the subspace \(\alpha_f(\mathbf{X} \times \{0\})\) is constant. There therefore exists a unique continuous map \(\varphi \colon \operatorname{Côn}(f) \to \mathbf{Z}\) such that \(h = \varphi \circ \pi\) , where \(\pi\) denotes the canonical surjection of \(\operatorname{Cyl}(f)\) onto \(\operatorname{Côn}(f)\) . The map \(\varphi\) satisfies \(\alpha = \varphi \circ \alpha_f'\) and \(\beta = \varphi \circ \beta_f'\) and is the only one having these properties, since the images of \(\alpha_f'\) and \(\beta_f'\) cover \(\operatorname{Côn}(f)\) .

Example 1. --- Let X be a topological space. We denote by C(X), and call the cone of the space X, the cone of the map IdX; it is the space obtained from X × I by contracting X × \{0\}. Let \(\pi\) be the canonical map from X × I to C(X); the image C'(X) of X × {[}0, 1{[} under \(\pi\) is called the open cone of the space X.

Suppose that X is not empty. Then the cone C(X) is not empty and its base point is still called the vertex of the cone.

The map \(((x,t),u)\mapsto(x,t(1-u))\) from \((\mathrm{X}\times\mathbf{I})\times\mathbf{I}\) into \(X\times I\) is a homotopy connecting the identity map of \(X\times I\) to the map \((x,t)\mapsto(x,0)\). From this we deduce, by passing to quotient sets, a map \(\sigma\colon\mathrm{C}(\mathrm{X})\times\mathbf{I}\to\mathrm{C}(\mathrm{X})\) such that \(\sigma(\pi(x,t),u)=\pi(x,t(1-u))\) for \(x\in X\), \(t,u\in I\). The map \(\sigma\) is continuous by I, p. 19, prop. 10. The map \(\sigma\) is then a pointed homotopy at s connecting the identity map of \(\mathrm{C}(\mathrm{X})\) to the constant map with image \(\{s\}\). Consequently, the cone \(\mathrm{C}(\mathrm{X})\) is contractible at its vertex s. Since \(\sigma(\mathrm{C}'(\mathrm{X})\times\mathbf{I})\) is contained in \(\mathrm{C}'(\mathrm{X})\), the map \(\sigma\) defines, by passing to subspaces, a pointed homotopy at s connecting the identity map of \(\mathrm{C}'(\mathrm{X})\) to the constant map with image \(\{s\}\); the open cone \(\mathrm{C}'(\mathrm{X})\) is therefore contractible at s.

Remarks. --- 1) Let X be a topological space and let A be a subspace of X; denote by i the canonical injection of A into X. The canonical retraction \(r_{i}\) of the cylinder \(\mathrm{Cyl}(i)\) onto its base maps the subspace \(\alpha_{i}(\mathrm{A} \times \{0\})\) into A. Let \(\rho: \mathrm{C}\hat{\mathrm{o}}\mathrm{n}(i) \to \mathrm{X}/\mathrm{A}\) be the continuous map induced therefrom by passing to quotients. Suppose that the pair (X, A) has the homotopy extension property. By III, p. 243, Theorem 1, the map \(r_{i}\) defines a homotopy equivalence of the pair \((\mathrm{Cyl}(i), \alpha_{i}(\mathrm{A} \times \{0\}))\) onto the pair (X, A). By Remark 4 of III, p. 253, the map \(\rho\) is then a homotopy equivalence of the pair \((\mathrm{C}\hat{\mathrm{o}}\mathrm{n}(i), s)\) onto the pair \((\mathrm{X}, s_{\mathrm{X}/\mathrm{A}})\). It is in particular a homotopy equivalence.

\begin{enumerate}
\def\labelenumi{\arabic{enumi})}
\setcounter{enumi}{1}
\tightlist
\item
  Let X and Y be topological spaces, and let \(f\colon X \to Y\) be a continuous map. The canonical map \(\alpha_{f}^{\prime}\colon X \times I \to \operatorname{Côn}(f)\) is constant on \(X \times \{0\}\) and therefore defines a continuous map \(\gamma_{f}\colon C(X) \to \operatorname{Côn}(f)\). The restriction to \(C'(X)\) of the map \(\gamma_{f}\) is injective and strict, and induces by passing to subspaces a homeomorphism of the open cone \(C'(X)\) onto the complement of Y in \(\operatorname{Côn}(f)\).
\end{enumerate}

Identify the base of the cone \(C(X)\) with the space X, and denote by \(u\) the continuous map from \(Y \cup_f C(X)\) into \(\operatorname{Côn}(f)\) induced by the maps \(\beta_f':Y\to\operatorname{Côn}(f)\) and \(\gamma_f:C(X)\to\operatorname{Côn}(f)\). Conversely, let \(v:\operatorname{Côn}(f)\to Y\cup_f C(X)\) be the unique continuous map such that \(v\circ\alpha_f'\) is the canonical map from \(X\times I\) onto \(C(X)\), and \(v\circ\beta_f'\) is the canonical map from \(Y\) to \(Y\cup_f C(X)\). The maps \(u\) and \(v\) are mutually inverse homeomorphisms.

